% Solutions to the Chapter 6 exercises, from lectures/L06/appendix/solutions: the README's prose
% and its answers to the reflection questions, each placed with the set it belongs to, and every
% source file typeset straight from the repository.

\section{Chapter 6: Multithreading and Synchronization}
\label{sol:sec:c6}
Reference implementations for the exercises in \secref{c6:sec:exercises}. Each program is a
directory under \file{lectures/L06/appendix/solutions}, built and run with \code{make}; its
Makefile is the one from \secref{c1:sec:makefile}, with \code{-pthread} added to the compiler
flags. Where a set's exercises build on one another, the solution is the program as it stands
after the last of them.

% ------------------------------------------------------------------------------------------
\subsection{Exercise Set 1: Threads and atomic stop flag}
\label{sol:set:c6.1}
Exercises~\ref{c6:ex:1.1} to \ref{c6:ex:1.3}, in \file{exercise1}.

\code{workerThread()} takes a print speed and an atomic stop flag. It prints a message at the given
interval until the stop flag is set. \code{main()} starts the thread, sleeps for a timeout, sets the
flag, and joins.

\cppfile{lectures/L06/appendix/solutions/exercise1/exercise1.cpp}

\needspace{8\baselineskip}
\subsubsection{Reflection: Exercise~\ref{c6:ex:1.2}}
\task{What happens if \code{join()} is not called?}
\begin{itemize}
  \item The \code{std::thread} object is destroyed while still joinable, which calls
    \code{std::terminate()} and crashes the program.
  \item The worker thread may not finish executing before the process exits.
\end{itemize}

\task{Which thread prints first, the main thread or the worker thread?}
\begin{itemize}
  \item Neither is guaranteed to print first.
  \item Both threads are ready to execute as soon as the worker thread is created, and the OS
    scheduler decides which runs first.
\end{itemize}

\task{Is the print order guaranteed?}
\begin{itemize}
  \item No. Thread scheduling is non-deterministic, so the order may differ between runs and
    between platforms.
\end{itemize}

\needspace{8\baselineskip}
\subsubsection{Reflection: Exercise~\ref{c6:ex:1.3}}
\task{Why is \code{std::atomic<bool>} suitable for this stop flag?}
\begin{itemize}
  \item It guarantees that reads and writes are indivisible and visible across threads without a
    data race.
  \item No additional synchronization primitive is needed for a single boolean flag that is written
    by one thread and read by another.
\end{itemize}

\task{Would a normal \code{bool} be safe here?}
\begin{itemize}
  \item No. Concurrent access to a plain \code{bool} from multiple threads without synchronization
    is undefined behavior.
  \item The compiler or CPU may also cache the value in a register, causing the worker thread to
    never see the updated value.
\end{itemize}

\task{Why is \code{load()} used when reading the flag?}
\begin{itemize}
  \item \code{load()} makes the atomic read explicit and ensures the value is fetched from the
    atomic object with the correct memory ordering, not from a stale cached copy.
\end{itemize}

\task{Why is \code{store()} used when writing the flag?}
\begin{itemize}
  \item \code{store()} makes the atomic write explicit and ensures the new value is published to all
    threads with the correct memory ordering.
\end{itemize}

% ------------------------------------------------------------------------------------------
\subsection{Exercise Set 2: Shared counter with data race}
\label{sol:set:c6.2}
Exercises~\ref{c6:ex:2.1} to \ref{c6:ex:2.3}, in \file{exercise2}.

\code{incrementCounter()} takes a counter and an iteration count. Two threads increment the same
counter concurrently. The exercises progress through three versions:
\begin{itemize}
  \item A plain \code{std::uint32_t} that demonstrates the data race.
  \item A mutex-protected version that fixes it.
  \item An \code{std::atomic<std::uint32_t>} that achieves the same result without a mutex.
\end{itemize}
The solution is the last of the three.

\cppfile{lectures/L06/appendix/solutions/exercise2/exercise2.cpp}

\needspace{8\baselineskip}
\subsubsection{Reflection: Exercise~\ref{c6:ex:2.1}}
\task{What final value do you expect?}
\begin{itemize}
  \item \code{200000}, since each of the two threads increments the counter \code{100000} times.
\end{itemize}

\task{Do you always get the expected value?}
\begin{itemize}
  \item No. Due to the data race the result is typically lower than \code{200000} and varies
    between runs.
\end{itemize}

\task{Why is this program unsafe?}
\begin{itemize}
  \item Two threads read, modify, and write the same counter without any synchronization.
  \item The read-modify-write sequence is not atomic for a plain integer, so one thread can
    overwrite the other's increment.
\end{itemize}

\task{Which condition makes this a data race?}
\begin{itemize}
  \item Two or more threads access the same memory location concurrently, at least one access is a
    write, and there is no synchronization between them.
\end{itemize}

\needspace{8\baselineskip}
\subsubsection{Reflection: Exercise~\ref{c6:ex:2.2}}
\task{What final counter value do you get now?}
\begin{itemize}
  \item Always \code{200000}.
\end{itemize}

\task{Why does the mutex solve the problem?}
\begin{itemize}
  \item The mutex ensures that only one thread can be inside the critical section at a time, making
    the read-modify-write sequence effectively atomic.
\end{itemize}

\task{Why is \code{std::lock_guard} safer than manually calling \code{lock()} and \code{unlock()}?}
\begin{itemize}
  \item \code{std::lock_guard} releases the lock automatically when it goes out of scope, even if an
    exception is thrown.
  \item Manual \code{lock()}/\code{unlock()} risks forgetting to unlock, or leaving the mutex locked
    if an exception occurs between the two calls.
\end{itemize}

\needspace{8\baselineskip}
\subsubsection{Reflection: Exercise~\ref{c6:ex:2.3}}
\task{What final counter value do you get?}
\begin{itemize}
  \item Always \code{200000}.
\end{itemize}

\task{Why is \code{std::atomic<std::uint32_t>} sufficient in this exercise?}
\begin{itemize}
  \item Incrementing a single integer is one independent operation. \code{std::atomic} makes the
    increment indivisible and thread-safe without needing a mutex.
\end{itemize}

\task{Would an atomic counter still be sufficient if several related variables had to be updated
together?}
\begin{itemize}
  \item No. If multiple related variables must be updated as a consistent unit, an atomic type
    alone cannot protect the compound operation.
  \item A mutex is needed to ensure the update of all related variables appears atomic to other
    threads.
\end{itemize}

% ------------------------------------------------------------------------------------------
\subsection{Exercise Set 3: Shared data protected by mutex}
\label{sol:set:c6.3}
Exercises~\ref{c6:ex:3.1} to \ref{c6:ex:3.4}, in \file{exercise3}.

\code{SharedMem} holds a \code{std::uint16_t} data value and a \code{newData} flag.
\code{txThread()} produces incrementing values and sets \code{newData}; \code{rxThread()} polls for
new data and consumes it. Both threads share one mutex that protects all accesses to
\code{SharedMem}. \code{main()} runs both threads for a fixed timeout, then sets the stop flag and
joins.

\cppfile{lectures/L06/appendix/solutions/exercise3/exercise3.cpp}

\needspace{8\baselineskip}
\subsubsection{Reflection: Exercise~\ref{c6:ex:3.4}}
\task{Why must \code{data} and \code{newData} be protected by the same mutex?}
\begin{itemize}
  \item They are logically related --- RX must read \code{data} and \code{newData} together as a
    consistent snapshot.
  \item If they were protected separately, TX could update \code{data} after RX has already read
    \code{newData}, causing RX to act on a stale or mismatched value.
\end{itemize}

\task{Why is \code{newData} not atomic in this design?}
\begin{itemize}
  \item \code{newData} must be checked and updated together with \code{data} as a consistent unit.
  \item Protecting \code{newData} with an atomic alone would not prevent TX from updating
    \code{data} between RX reading the flag and reading the value.
  \item The mutex protects both fields together, making the compound operation safe.
\end{itemize}

\task{What could happen if TX updated \code{data} and \code{newData} without locking the mutex?}
\begin{itemize}
  \item RX could read \code{newData = true} before \code{data} has been written, or read a
    partially updated \code{data} value.
  \item This is a data race and results in undefined behavior.
\end{itemize}

\task{Why is the stop flag atomic instead of mutex-protected?}
\begin{itemize}
  \item The stop flag is a single boolean written by one thread and read by others. No compound
    operation is needed.
  \item An atomic boolean is sufficient and avoids the overhead of locking the mutex on every loop
    iteration just to check a flag.
\end{itemize}

\task{Could \code{std::atomic<SharedMem>} be used instead of a mutex?}
\begin{itemize}
  \item Not without additional complexity. RX's operation is a read-modify-write sequence: load the
    struct, check \code{newData}, clear \code{newData}, and store back.
  \item Even with \code{std::atomic<SharedMem>}, that sequence is not atomic as a whole. TX could
    store a new value between RX's load and RX's store, causing RX to silently overwrite TX's new
    data.
  \item Making it correct without a mutex would require a compare-exchange (CAS) retry loop in RX,
    which is significantly more complex than the mutex approach.
  \item The mutex is the right tool here because it makes the entire compound operation safe
    through mutual exclusion.
\end{itemize}

% ------------------------------------------------------------------------------------------
\subsection{Exercise Sets 4 and 5: Thread-safe stub driver}
\label{sol:set:c6.4}\label{sol:set:c6.5}
Exercises~\ref{c6:ex:4.1} to \ref{c6:ex:5.3}, in \file{exercise4-5}.

\code{driver::counter::Interface} defines a pure virtual interface with \code{isInitialized()},
\code{value()}, \code{increment()}, and \code{reset()}.

\code{driver::counter::Stub} implements the interface:
\begin{itemize}
  \item \code{myMutex} (\code{mutable std::mutex}) protects \code{myValue}.
  \item \code{myInitialized} (\code{std::atomic<bool>}) guards all counter operations --- each method
    returns early if the driver is not initialized.
  \item Copy and move are deleted.
\end{itemize}
\code{main()} creates one \code{Stub}, starts two threads via \code{counterThread()} (100 and 200
increments respectively), joins them, and prints the final value --- expected \code{300}.

\cppfile{lectures/L06/appendix/solutions/exercise4-5/include/driver/counter/interface.hpp}
\cppfile{lectures/L06/appendix/solutions/exercise4-5/include/driver/counter/stub.hpp}
\cppfile{lectures/L06/appendix/solutions/exercise4-5/source/main.cpp}

\needspace{8\baselineskip}
\subsubsection{Reflection: Exercise~\ref{c6:ex:4.3}}
\task{Why does \code{value()} need to lock the mutex even though it only reads the value?}
\begin{itemize}
  \item Another thread could be calling \code{increment()} or \code{reset()} concurrently.
  \item Reading \code{myValue} while it is being written is a data race. The mutex ensures the read
    is not interleaved with a write.
\end{itemize}

\task{Why is \code{myMutex} marked \code{mutable}?}
\begin{itemize}
  \item \code{value()} is a \code{const} method and cannot modify member variables.
  \item Locking a mutex is an implementation detail, not a logical mutation of the object.
  \item \code{mutable} allows the mutex to be locked inside a \code{const} method without violating
    \code{const} correctness.
\end{itemize}

\task{What would happen if \code{increment()} did not lock the mutex?}
\begin{itemize}
  \item Two threads could read the same value, both increment it, and write back the same result,
    losing one increment.
  \item The counter would end up lower than expected --- identical to the data race in
    Exercise~\ref{c6:ex:2.1}.
\end{itemize}

\task{Could \code{std::atomic<std::uint32_t>} be used instead in this specific class?}
\begin{itemize}
  \item Yes. In this specific class each operation (increment, read, reset) acts on a single
    variable independently, so \code{std::atomic<std::uint32_t>} would be sufficient.
\end{itemize}

\task{Why might a mutex still be preferable in a more complex driver?}
\begin{itemize}
  \item If future requirements add more member variables that must be updated together with
    \code{myValue}, a mutex can protect all of them as a unit.
  \item An atomic can only protect a single variable; compound multi-variable operations still
    require a mutex.
\end{itemize}

\needspace{8\baselineskip}
\subsubsection{Reflection: Exercise~\ref{c6:ex:5.3}}
\task{Why is \code{myInitialized} suitable as an atomic variable?}
\begin{itemize}
  \item It is a single boolean flag that is read and written independently.
  \item No other variable needs to be updated in the same atomic step, so
    \code{std::atomic<bool>} is sufficient.
\end{itemize}

\task{Why is \code{myValue} still protected by a mutex?}
\begin{itemize}
  \item \code{myValue} is accessed by multiple operations (\code{increment}, \code{reset},
    \code{value}) that must not interleave.
  \item A mutex protects all operations on \code{myValue} as a group, preventing data races.
\end{itemize}

\task{Would it be safe to make both variables atomic if future requirements added more related
state?}
\begin{itemize}
  \item No. If future state must be updated together with \code{myValue} as a consistent unit,
    atomic types alone cannot protect the compound update.
  \item A mutex would be needed to keep the update of all related variables atomic as a group.
\end{itemize}

\task{What is the main design difference between a thread-safe flag and thread-safe shared data?}
\begin{itemize}
  \item A thread-safe flag is a single boolean that only needs to be observed independently --- an
    atomic is sufficient.
  \item Thread-safe shared data involves multiple related fields that must be read or written
    together as a consistent snapshot, requiring a mutex to protect the compound operation.
\end{itemize}

\task{What could happen if \code{setInitialized(false)} is called between the initialization check
and the mutex lock in \code{increment()}? How would you fix this?}
\begin{itemize}
  \item \code{increment()} would pass the initialization check, then the driver would be marked
    uninitialized by another thread, and \code{increment()} would proceed to modify \code{myValue}
    --- violating the intended guard.
  \item The fix is to make the check and the update one critical section: check
    \code{myInitialized} after taking the lock in \code{increment()} (and likewise in
    \code{value()} and \code{reset()}), \textbf{and} have \code{setInitialized()} take the same
    mutex when it changes the flag.
  \item Moving the check inside the lock is not enough on its own: \code{setInitialized()} does not
    lock the mutex, so it could still change the flag after \code{increment()} has checked it. With
    both sides under the mutex, \code{setInitialized(false)} either completes before
    \code{increment()} checks the flag, or waits until the increment is done.
\end{itemize}

% ------------------------------------------------------------------------------------------
\subsection{Exercise Set 6: Condition variable}
\label{sol:set:c6.6}
Exercise~\ref{c6:ex:6.1}, in \file{exercise6}.

Built on Exercise Set 3. \code{rxThread()} no longer polls with a sleep; instead it acquires a
\code{std::unique_lock} and calls \code{cv.wait()} with \code{hasNewData()} as the predicate.
\code{txThread()} calls \code{cv.notify_one()} after releasing the lock. \code{main()} sets the stop
flag while holding the mutex, so that it cannot change between the receiver checking the predicate
and going to sleep, and then calls \code{cv.notify_all()} so the receiver wakes up and exits
cleanly.

\cppfile{lectures/L06/appendix/solutions/exercise6/exercise6.cpp}

\needspace{8\baselineskip}
\subsubsection{Reflection: Exercise~\ref{c6:ex:6.1}}
\task{Why is \code{std::unique_lock} required here instead of \code{std::lock_guard}?}
\begin{itemize}
  \item \code{cv.wait()} must unlock the mutex while the thread is sleeping and re-lock it when
    woken up.
  \item \code{std::lock_guard} has no way to release and re-acquire the lock --- it only supports
    locking on construction and unlocking on destruction.
  \item \code{std::unique_lock} provides \code{lock()} and \code{unlock()}, which \code{cv.wait()}
    uses internally.
\end{itemize}

\task{What is a spurious wake-up, and how does the predicate protect against it?}
\begin{itemize}
  \item A spurious wake-up is when \code{cv.wait()} returns even though no \code{notify_one()} or
    \code{notify_all()} was called, due to OS implementation details.
  \item The predicate re-checks the condition after every wake-up. If the condition is not met,
    \code{wait()} automatically re-enters the waiting state instead of proceeding.
\end{itemize}

\task{Why must the stop flag also be checked inside the predicate?}
\begin{itemize}
  \item When the stop flag is set and \code{notify_all()} is called, the receiver wakes up.
  \item Without \code{stop.load()} in the predicate, the receiver would see
    \code{hasNewData() == false}, treat it as a spurious wake-up, and go back to sleep --- never
    exiting cleanly.
\end{itemize}

\task{How does this design differ from the polling approach in Exercise~\ref{c6:ex:3.3}?}
\begin{itemize}
  \item In Exercise~\ref{c6:ex:3.3} the receiver wakes every \code{100 ms} regardless of whether
    data is available, wasting CPU cycles and adding up to \code{100 ms} of latency.
  \item With a condition variable the receiver sleeps until TX explicitly wakes it, reducing
    latency to near zero and eliminating unnecessary polling.
\end{itemize}

\task{What would happen if \code{notify_all()} was not called after setting the stop flag?}
\begin{itemize}
  \item The receiver thread would remain blocked on \code{cv.wait()} even after \code{stop} is set
    to \code{true}, because nothing wakes it up.
  \item The main thread would then block forever on \code{join()}, and the program would hang.
\end{itemize}

\task{Why is \code{notify_all()} used here instead of \code{notify_one()}?}
\begin{itemize}
  \item In this design there is only one receiver thread, so \code{notify_one()} would be
    sufficient here.
  \item \code{notify_all()} is used for correctness in the general case --- if multiple threads were
    waiting on the same condition variable, \code{notify_one()} would only wake one of them,
    leaving the others blocked indefinitely.
\end{itemize}

% ------------------------------------------------------------------------------------------
\subsection{Exercise Set 7: \code{std::async} and \code{std::future}}
\label{sol:set:c6.7}
Exercises~\ref{c6:ex:7.1} and \ref{c6:ex:7.2}, in \file{exercise7}.

\code{validateFirmware()} simulates a slow flash read by sleeping for \code{2000 ms}, then iterates
over the buffer looking for \code{0xFFU}. It is launched asynchronously via
\code{std::async(std::launch::async, ...)}. \code{main()} polls
\code{future.wait_for(std::chrono::milliseconds(200U))} against \code{std::future_status::ready} in
a loop, printing a status message each iteration, then retrieves the result with
\code{future.get()}.

\cppfile{lectures/L06/appendix/solutions/exercise7/exercise7.cpp}

\needspace{8\baselineskip}
\subsubsection{Reflection: Exercise~\ref{c6:ex:7.1}}
\task{What does \code{std::launch::async} guarantee compared to omitting the launch policy?}
\begin{itemize}
  \item \code{std::launch::async} guarantees the function runs in a new thread immediately,
    independently of when \code{get()} is called.
  \item Without the policy, the implementation may choose \code{std::launch::deferred}, which
    defers execution until \code{get()} is called --- running the function synchronously on the
    calling thread instead.
\end{itemize}

\task{At what point does the main thread block?}
\begin{itemize}
  \item The main thread blocks when \code{future.get()} is called.
  \item Everything between \code{std::async} and \code{future.get()} executes concurrently while
    validation runs in the background.
\end{itemize}

\task{What would happen if \code{future.get()} was called immediately after \code{std::async}?}
\begin{itemize}
  \item The main thread would block immediately, waiting for \code{validateFirmware()} to complete
    before continuing.
  \item No concurrent work would happen --- it would behave like a synchronous function call.
\end{itemize}

\needspace{8\baselineskip}
\subsubsection{Reflection: Exercise~\ref{c6:ex:7.2}}
\task{How does \code{wait_for()} differ from \code{get()}?}
\begin{itemize}
  \item \code{wait_for()} waits at most for the given duration and returns a status
    (\code{ready}, \code{timeout}, or \code{deferred}) without consuming the result or blocking
    indefinitely.
  \item \code{get()} blocks until the result is available and then returns or re-throws it,
    consuming the future.
\end{itemize}

\task{Is it safe to call \code{get()} after \code{wait_for()} returns
\code{std::future_status::ready}?}
\begin{itemize}
  \item Yes. Once \code{wait_for()} returns \code{std::future_status::ready} the result is
    guaranteed to be available, and \code{get()} will return immediately without blocking.
\end{itemize}
